\section{Arithmetic, Dimension, and the Neural Verification Account}
\label{app:r41arithmetic}
This appendix supplies the arithmetic and work details for
Section~\ref{sec:r41direct} of the article. The trained functions at both
successive dates are inputs to the certificate. Approximation of those
functions by their former training targets is not an allowed operation.

\subsection{Exact activation regions and the terminal discrepancy}
Write the stored scalar network as
\[
 f(x)=b_0+\sum_{j=1}^m v_j(w_jx+b_j)_+.
\]
All coefficients are rational when read as the exact values of their binary
representations. Units with $w_j=0$ contribute constants. For the remaining
units, order and deduplicate the rational roots $-b_j/w_j$ in $[0,1]$.
Between consecutive roots the active set is constant, and summing its
coefficients gives an exact affine function $A_kx+B_k$. Negative slopes do
not require a separate approximation rule. The maximum of $|A_k|$ over
nonempty regions is an all-state Lipschitz bound. This maximum can be much
smaller than the sum of absolute weights.

\begin{proposition}[Evaluation after rational region compilation]
\label{prop:r41compile}
Suppose every rounded root differs from its exact root by at most $e_r$,
while order is retained. Let $L=\max_k|A_k|$. If region selection at $x$
uses the rounded roots, evaluating the selected exact affine extension
instead of $f(x)$ makes an error no larger than $2Le_r$. If several rounded
roots coincide, the same bound holds provided every displaced boundary
between the chosen and true regions lies within $e_r$ of $x$. Outward
coefficient evaluation at a midpoint $x_c$ of an input interval $I$ and an
additional $L\operatorname{rad}(I)$ then enclose $f(I)$, after charging
midpoint rounding as well.
\end{proposition}
\begin{proof}
An incorrectly selected region is separated from the true one only by
roots whose rounded comparisons differ at $x$. Such roots are within
$e_r$ of $x$. Choose the boundary of the selected region closest to $x$;
call it $r$. The selected affine function equals $f$ at $r$ by continuity.
Both $f$ and that affine function have Lipschitz constant at most $L$.
Their difference at $x$ is therefore at most $2L|x-r|\leq2Le_r$.
For an interval, add the Lipschitz variation from its midpoint and the
outward arithmetic enclosure of the selected affine evaluation.
\end{proof}

For the terminal primitive, split further at the cost kink $3/8$.
On every resulting interval $f_T-g$ is a quadratic polynomial with rational
coefficients. A continuous quadratic achieves its extrema at an endpoint
or at its stationary point if that point is in the interval. Testing those
points gives rational global endpoints. The terminal error is generally
nonzero. In the executed study, 66 candidates yield a terminal interval
approximately $[-0.000668000,0.000410383]$ for every economy, because the
terminal network and terminal primitive are shared.

\subsection{Rounding the specified execution graph}
A range-error pair consists of an exact range $[a,b]$ and a nonnegative
rational error $e$. It means that an ideal node $y$ lies in $[a,b]$ and its
implemented value differs by at most $e$. Put $M([a,b])=\max(|a|,|b|)$.
For separate round-to-nearest binary32 operations, use
$u=2^{-24}$ and $\alpha=2^{-150}$. If a sum has pre-rounding propagated
error $e_s=e_1+e_2$ and ideal range $I_s$, its new error is
\[
 e_s+u\{M(I_s)+e_s\}+\alpha.
\]
For a product, first put
\[
 e_p=M(I_1)e_2+M(I_2)e_1+e_1e_2;
\]
then add $u\{M(I_1I_2)+e_p\}+\alpha$. These formulae follow by expanding
$(y_1+\delta_1)(y_2+\delta_2)-y_1y_2$ and applying the elementary rounding
inequality. Rational range propagation is performed separately. Conversion
of the state, action, price, and shock is included at the input nodes.
Clipping to a representable closed interval is nonexpansive and cannot
increase this error.

The deposited implementation applies these recurrences to the operation
order used in the binary32 transition, cost, and terminal routines.
Its verification test compares 1,032 native executions with exact rational
execution of the corresponding ideal graphs. Those finite tests detect
implementation mistakes; the range induction, not the tests, supplies the
all-state guarantee. The maximum observed error is less than $0.224$ of
the corresponding uniform allowance. No conclusion is drawn for a fused
multiply-add implementation or a platform that flushes subnormals to zero.

Nearest-node policy selection deserves a separate argument because the
actor is discontinuous. For a true state $x$ and acquired state
$\widehat x$, a node selected nearest to $\widehat x$ is at most
$h_x+|\widehat x-x|$ from $x$. The chosen-action and Bellman-minimum
comparisons each pay $L_x|\widehat x-x|$. There is consequently no need
for an unjustified Lipschitz bound on the actor itself. At the dyadic grids
in this study, stored action values and state cutpoints are exactly
representable in binary32; that fact does not remove errors in state
acquisition, transition, or cost evaluation.

\subsection{Continuous innovations and multidimensional verification}
\label{app:r41work}
The direct certificate is not tied to scalar states or finite-support
uncertainty. This extension states explicitly what additional work is
required; it does not identify a grid verifier with a dimension-free
method. Let $X\subset\mathbb R^{d_x}$, $A\subset\mathbb R^{d_a}$, and
$Z\subset\mathbb R^{d_z}$ be compact. The innovation has a known law and
admits a deterministic quantizer $q$ with
$\|Z-q(Z)\|\leq h_z$ almost surely. The probabilities of the quantizer
cells are supplied exactly or by enclosing arithmetic. Take a law-invariant,
monotone, cash-invariant one-step risk functional, evaluated conditionally
on the current state and action. Expectation and entropic risk satisfy this
requirement. Suppose the transition is Lipschitz in the innovation with
constant $L_{Fz,t}$. If stage costs also depend on the innovation, add their
innovation Lipschitz contribution in the definition below.

\begin{theorem}[A dimension-explicit neural certificate with quantized innovations]
\label{thm:r41dimension}
Let the hypotheses of the direct neural certificate hold on $X$ and $A$.
Put $\zeta_t=\beta_t L_{t+1}L_{Fz,t}h_z$. Construct nodal intervals using
$q(Z)$ rather than $Z$. If $(\ell_t,u_t,\eta_t)$ are computed from these
intervals, then
\[
 (\ell_t-\zeta_t,\ u_t+\zeta_t,\ \eta_t+2\zeta_t)
\]
are valid residual and actor bounds for the unquantized economy. The global
policy certificate increases by at most $4\sum_{t<T}B_t\zeta_t$.

Let $K_x,K_a,m$ be the numbers of state nodes, action nodes, and quantizer
cells. Let $P_t$ count scalar arithmetic operations in a forward network
evaluation, $C_{F,t}$ the transition and cost evaluation work, and
$C_{\rho,t}(m,p)$ the work to enclose one risk evaluation at precision $p$.
In units of certified precision-$p$ arithmetic, a streaming direct
verification uses
\begin{align*}
 W\leq C\sum_{t<T}K_xK_a
 \{m(C_{F,t}+P_{t+1})+C_{\rho,t}(m,p)\}
 +C\sum_{t\leq T}K_xP_t+W_T+W_{\rm io}.
\end{align*}
Here $W_T$ certifies the terminal discrepancy and $W_{\rm io}$ includes
serialization and the durable evidence writes. For a batch of $b$ state
nodes, its working storage, in scalar words, is at most
\[
 O\!\left(\sum_{t\leq T}P_t+TK_xd_a+md_z+
 bK_am\,w_{\max}+K_x\right),
\]
where $w_{\max}$ bounds the largest network layer width and $b\leq K_x$.
The policy output itself requires $TK_xd_a$ words. Tensor covers of bounded
boxes can be chosen with
\[
 K_x=O(h_x^{-d_x}),\qquad K_a=O(h_a^{-d_a}),\qquad
 m=O(h_z^{-d_z}).
\]
The hidden constants depend on the domains and the norm, not on training
convergence. A finite refinement schedule is charged by summing these
bounds over every attempted resolution, not just the accepted resolution.
\end{theorem}
\begin{proof}
Couple $Z$ and $q(Z)$ on their defining probability space. The two future
network values differ almost surely by at most
$L_{t+1}L_{Fz,t}h_z$. Monotonicity and cash invariance bound the difference
of their risk evaluations by the same number. Law invariance identifies
the quantized risk evaluation with its finite-support calculation.
Multiplication by $\beta_t$ gives a uniform $Q$ discrepancy $\zeta_t$.
Taking an infimum does not increase a uniform error. A residual therefore
pays $\zeta_t$ at each endpoint; a chosen-action excess pays $2\zeta_t$.
The residual oscillation and actor allowance together pay $4\zeta_t$.
Backward induction gives the stated policy bound.

For each of the $TK_xK_a$ state-action pairs the program evaluates $m$
transitions and networks and then one risk aggregate. Current networks
are evaluated at the state nodes as well. Streaming over state batches
retains only the batch's forward activations, the nodal extrema and stored
actors, the innovation support, and the network coefficients. This proves
the operation and storage counts. An equally spaced tensor cover proves
the node bounds. Exact enumeration of a policy's innovation tree at $Q$
initial queries adds, separately, order
$Q\sum_{t=0}^Tm^t$ evaluations and a depth-first stack of order $T$ times
the stored state and risk-aggregation data. This exponential query cost is
not hidden inside the Bellman verification count.
\end{proof}

The notation $C_{\rho,t}(m,p)$ is intentional. A native addition is not an
exact exponential, logarithm, or certified inverse. Bit costs require a
specified arithmetic implementation and must multiply the appropriate
operation counts by its precision-dependent costs. The complete-step
matrix result in the article similarly exposes its lower spectral bound
and residual enclosure rather than suppressing conditioning in an
unspecified constant. Neural Lipschitz bounds, risk-evaluation ranges,
terminal verification, and the admissible condition margins must be
reported before either count is used numerically.

A refinement reduces verification slack; it cannot force a fixed inaccurate
network to satisfy an arbitrary policy target. More precisely, suppose
current-value and quantized $Q$ nodal intervals have widths at most
$w_{f,t}$ and $w_{q,t}$. Choosing a minimum-upper-endpoint action gives the
following useful upper bound on the certificate:
\begin{align*}
 E_f\leq &\sum_{t<T}B_t\operatorname{osc}(f_t-T_tf_{t+1})
       +B_T\operatorname{osc}(f_T-g)\\
 &+\sum_{t<T}B_t\{(2L_t+4L_{x,t})h_x+3L_{a,t}h_a
              +4\zeta_t+2w_{f,t}+3w_{q,t}\}.
\end{align*}
To check the coefficients, the grid minimum differs from the continuous
minimum by between zero and $L_{a,t}h_a$. The residual range therefore
pays at most $2L_{a,t}h_a+2w_{f,t}+2w_{q,t}$, in addition to its state
cover. The selected action pays at most one further $L_{a,t}h_a+w_{q,t}$
and its two state movements. Quantization supplies the four allowances
already proved. Thus the refinement tends to the network's residual and
terminal oscillations, not automatically to zero.

For prescribed verification slack, this theorem specifies its dependence
on state, action and innovation dimension, horizon, network size, precision,
and the problem's Lipschitz and conditioning data. It does not establish a
competitive high-dimensional work frontier. The new executed nonlinear
certificate below remains scalar; the higher dimensions in the quadratic
study do not change that fact. A coupled nonlinear scaling comparison with
strong sparse, projection, or fitted dynamic-programming baselines remains
an empirical question distinct from this extension theorem.
